\documentclass[11pt]{article}
\usepackage[a4paper,margin=1in]{geometry}
\usepackage[T1]{fontenc}
\usepackage{lmodern}
\usepackage{amsmath,amssymb,mathtools}
\usepackage{booktabs,array,longtable}
\usepackage{microtype}
\usepackage[hidelinks]{hyperref}
\usepackage{enumitem}
\usepackage{siunitx}
\usepackage{fancyhdr}
\usepackage{xcolor}
\usepackage{caption}
\setlength{\parindent}{1.2em}
\setlength{\parskip}{0.25em}
\setlength{\headheight}{14pt}
\pagestyle{fancy}
\fancyhf{}
\lhead{A Ternary $(4k-1(2))/3$ Collatz-Type Map}
\rhead{\thepage}
\newcommand{\vthree}{\nu_3}
\newcommand{\Sset}{\mathcal S}
\newcommand{\Tmap}{T}
\newcommand{\Mmap}{M}
\newcommand{\R}{r}

\title{\textbf{A Ternary $(4k-1(2))/3$ Collatz-Type Map}\\[4pt]
\large Complete 3-Adic Reduction, Algebraic Cycle Analysis, and Exhaustive Computational Verification up to $10^{11}$}
\author{\textbf{Farhad Banazadeh}\\
\href{mailto:fn.bana@hotmail.com}{fn.bana@hotmail.com}\\
\href{https://orcid.org/0009-0004-7023-0298}{ORCID: 0009-0004-7023-0298}}
\date{18 September 2026}

\begin{document}
\maketitle

\begin{abstract}\sloppy
We study a ternary Collatz-type dynamical system on the positive integers not divisible by $3$.  In the compact notation of the title, $(4k-1(2))/3$ means that the affine numerator subtracts $1$ or $2$ according to the residue class of the current state modulo $3$; the actual accelerated map removes \emph{all} factors of $3$ at each step.  More precisely, for
\[
\Sset=\{x\in\mathbb Z_{>0}:3\nmid x\},\qquad r(x)=x\bmod 3\in\{1,2\},
\]
we define
\[
\Mmap(x)=4x-r(x),\qquad
\Tmap(x)=\frac{\Mmap(x)}{3^{\vthree(\Mmap(x))}}.
\]
The map has the fixed points $1$ and $2$ and the positive four-cycle
$22\to29\to38\to50\to22$.  We derive exact residue branches and inverse branches, obtain a closed Diophantine equation for a general cycle, and show how the known positive and negative cycles are reconstructed algebraically.  The four-cycle is governed by the near equality $4^4=256$ and $3^5=243$, giving the small cycle denominator $4^4-3^5=13$.

A geometric valuation heuristic gives
$\Pr(\vthree=k)=2/3^k$ for $k\ge1$, hence
$\mathbb E[\vthree]=3/2$ and the average logarithmic drift
$\log 4-\frac32\log 3\approx-0.261624$, with geometric factor about $0.7698$.  This is a probabilistic explanation for typical contraction, not a proof of global convergence.

Finally, we combine the previously archived exhaustive scan through $10^9$ with a new disjoint V2 scan of $10^9<x\le10^{11}$.  The resulting verified domain contains exactly $66,666,666,667$ admissible starts.  The fixed point $1$ attracts $6,602,925,203$ starts ($9.904388\%$), the fixed point $2$ attracts $43,211,124,603$ ($64.816687\%$), and the four-cycle attracts $16,852,616,861$ ($25.278925\%$).  No additional positive cycle, unresolved trajectory, or arithmetic overflow is reported in either component scan.  Across the complete verified domain, the trajectories contain exactly $5,577,182,144,312$ reduced-map iterations and remove $8,429,151,431,804$ factors of $3$ in total.  The full-domain record transient is $574$ reduced-map iterations from $97,312,923,676$; the record accumulated $3$-adic valuation is $744$ from the same start; and the largest fully reduced state observed is $580,315,803,007,928,858,285$, reached by a trajectory beginning at $52,991,345,843$.  The result is an exhaustive finite verification through $10^{11}$, not a theorem for all positive integers.
\end{abstract}

\noindent\textbf{Keywords:} Collatz-type map; ternary dynamics; 3-adic valuation; accelerated iteration; Diophantine cycles; exhaustive computation.

\section{Introduction}
The classical Collatz problem is one of the simplest-looking and most persistent open problems in discrete mathematics.  Its difficulty comes from the interaction between an affine growth rule and residue-dependent division.  Generalized Collatz systems replace parity by other moduli, use different affine coefficients, or compress several divisions into a single dynamical step.  In these systems, arithmetic residue classes can determine both the typical drift and the possible periodic structures.

The present map is a ternary accelerated system.  Every state is a positive integer not divisible by $3$.  The state is multiplied by $4$, its residue modulo $3$ is subtracted so that the result becomes divisible by $3$, and then every available factor of $3$ is removed.  The complete removal of powers of $3$ is the key structural feature: unlike an unaccelerated map that divides by only one factor of $3$, one reduced iteration may remove several factors at once.

This distinction also creates several natural computational statistics.  A trajectory has a number of reduced-map iterations, a total number of individual divisions by $3$, and several possible notions of peak: the maximum fully reduced state, the maximum value after only the compulsory first division by $3$, and the maximum raw affine value $\Mmap(x)$.  These quantities should not be conflated.

The paper has four main aims.  First, we formulate the map and its two residue branches precisely.  Second, we analyze the valuation distribution and the associated logarithmic contraction heuristic.  Third, we derive the exact cycle equation and reconstruct the known cycles algebraically.  Fourth, building on the exhaustive $10^9$ audit in the base article~\cite{banazadeh2026base}, we extend the verified positive domain to $10^{11}$ by combining that complete base interval with a disjoint V2 scan of $10^9<x\le10^{11}$.

\section{Definition of the domain and map}
Let
\begin{equation}
\Sset=\{x\in\mathbb Z_{>0}:x\not\equiv0\pmod3\}.
\end{equation}
Every $x\in\Sset$ has exactly one residue
\[
r(x)\in\{1,2\},\qquad r(x)\equiv x\pmod3.
\]
Define the affine pre-reduction map
\begin{equation}
\Mmap(x)=4x-r(x)=
\begin{cases}
4x-1,&x\equiv1\pmod3,\\
4x-2,&x\equiv2\pmod3.
\end{cases}
\label{eq:M}
\end{equation}
Since $4\equiv1\pmod3$, we have $4x-r(x)\equiv0\pmod3$.  Therefore $\Mmap(x)$ always contains at least one factor of $3$.

For $n\ne0$, let
\[
\vthree(n)=\max\{j\ge0:3^j\mid n\}
\]
be the $3$-adic valuation.  The accelerated map is
\begin{equation}
\boxed{
\Tmap(x)=\frac{4x-r(x)}{3^{\vthree(4x-r(x))}}
}
\label{eq:T}
\end{equation}
and is well-defined from $\Sset$ to itself.

\subsection{The title notation $(4k-1(2))/3$}
The compact notation in the title records the two ternary affine branches: subtract $1$ when the state is $1\pmod3$ and subtract $2$ when it is $2\pmod3$.  Equation~\eqref{eq:T} is the precise definition used throughout the mathematics and computation.  The denominator $3$ in the compact title emphasizes the compulsory ternary division, while the actual reduced map divides by the full power $3^{\vthree}$.

\section{Algebra of the two residue branches}
The branch structure becomes especially transparent after writing the state in residue form.

If $x=3q+1$, then
\begin{equation}
\Mmap(x)=4(3q+1)-1=12q+3=3(4q+1),
\end{equation}
so
\begin{equation}
\Tmap(3q+1)=\frac{4q+1}{3^{\vthree(4q+1)}}.
\label{eq:branch1}
\end{equation}
Likewise, if $x=3q+2$, then
\begin{equation}
\Mmap(x)=4(3q+2)-2=12q+6=3(4q+2),
\end{equation}
and therefore
\begin{equation}
\Tmap(3q+2)=\frac{4q+2}{3^{\vthree(4q+2)}}.
\label{eq:branch2}
\end{equation}
Equations~\eqref{eq:branch1}--\eqref{eq:branch2} show explicitly that one factor of $3$ is compulsory in the affine form, while any additional factors depend on the smaller branch quantities $4q+1$ and $4q+2$.

\subsection{Exact inverse branches}
The one-step relation can be written
\begin{equation}
3^k y=4x-r,
\qquad r\in\{1,2\},\quad k\ge1,
\label{eq:onestep}
\end{equation}
where $y=\Tmap(x)$.  Hence every possible predecessor associated with a prescribed residue $r$ and valuation $k$ has the form
\begin{equation}
x=\frac{3^k y+r}{4}.
\label{eq:inverse}
\end{equation}
Such a predecessor exists exactly when the numerator is divisible by $4$ and its valuation is exactly compatible with the chosen $k$.  Because $4\equiv1\pmod3$, a valid integer in \eqref{eq:inverse} automatically has residue $r$ modulo $3$.  This inverse-branch representation is useful both for cycle algebra and for constructing the preimage tree of an attractor.

\section{Step-by-step numerical examples}
The smallest states already reveal all three positive attractors.

For $x=1$,
\[
\Mmap(1)=4-1=3,\qquad \vthree(3)=1,\qquad \Tmap(1)=1,
\]
so $1$ is a fixed point.  For $x=2$,
\[
\Mmap(2)=8-2=6,\qquad \vthree(6)=1,\qquad \Tmap(2)=2,
\]
so $2$ is also a fixed point.

For $x=4$,
\[
\Mmap(4)=16-1=15=3\cdot5,
\qquad \Tmap(4)=5.
\]
For $x=5$,
\[
\Mmap(5)=20-2=18=3^2\cdot2,
\qquad \Tmap(5)=2,
\]
so $4\to5\to2$.  For $x=7$,
\[
\Mmap(7)=28-1=27=3^3,
\qquad \Tmap(7)=1,
\]
showing how a high valuation can produce a sharp fall in a single reduced step.

The nontrivial positive cycle is
\begin{equation}
22\longrightarrow29\longrightarrow38\longrightarrow50\longrightarrow22.
\label{eq:fourcycle}
\end{equation}
Indeed,
\[
\begin{aligned}
4(22)-1&=87=3\cdot29,\\
4(29)-2&=114=3\cdot38,\\
4(38)-2&=150=3\cdot50,\\
4(50)-2&=198=3^2\cdot22.
\end{aligned}
\]
Thus three steps remove one factor of $3$, while the final step removes two.  The valuation vector is $(1,1,1,2)$ and its sum is $K=5$.

\section{The 3-adic valuation heuristic}
The large-scale behavior is strongly influenced by the number of powers of $3$ removed in each reduced iteration.  Suppose, heuristically, that after the compulsory divisibility by $3$, the remaining residue information is distributed uniformly modulo higher powers of $3$.  Then for $k\ge1$,
\begin{equation}
\Pr\bigl(\vthree(\Mmap(x))=k\bigr)
=\left(\frac13\right)^{k-1}\left(1-\frac13\right)
=\frac{2}{3^k}.
\label{eq:geom}
\end{equation}
This is a geometric distribution on $\{1,2,3,\ldots\}$.

\subsection{Expected valuation}
From \eqref{eq:geom},
\begin{equation}
\mathbb E[\vthree]
=\sum_{k=1}^{\infty} k\frac{2}{3^k}
=\frac32.
\label{eq:meanv}
\end{equation}
Thus the heuristic average denominator in multiplicative scale is $3^{3/2}=3\sqrt3\approx5.196$, while the numerator has leading coefficient $4$.

\subsection{Average logarithmic drift}
For large $x$, the subtraction of $r(x)\in\{1,2\}$ is lower order, so
\[
\log\frac{\Tmap(x)}{x}
\approx \log4-\vthree\log3.
\]
Taking the heuristic expectation gives
\begin{equation}
\Delta_{\mathrm{heur}}
=\log4-\frac32\log3
\approx-0.261624.
\label{eq:drift}
\end{equation}
The associated geometric factor is
\begin{equation}
\lambda_{\mathrm{heur}}
=e^{\Delta_{\mathrm{heur}}}
=\frac{4}{3^{3/2}}
=\frac{4}{3\sqrt3}
\approx0.769800.
\label{eq:lambda}
\end{equation}
In this model a typical logarithmic step contracts the scale by about $23\%$.

\subsection{What the heuristic does not prove}
The negative drift gives a useful statistical explanation for the computational behavior, but it is not a proof that every deterministic orbit converges.  Successive valuations along a fixed orbit need not behave as independent samples from \eqref{eq:geom}, and a negative mean drift does not by itself exclude exceptional unbounded trajectories.  The exhaustive results below are therefore stated only for the finite domain actually tested.

\section{Algebraic cycle equation}
This section gives the exact arithmetic structure behind every periodic orbit of the map.

Suppose a cycle has length $L$ and states
\[
x_0,x_1,\ldots,x_{L-1},\qquad x_L=x_0.
\]
At step $i$, define
\[
r_i=x_i\bmod3\in\{1,2\},
\qquad
k_i=\vthree(4x_i-r_i)\ge1.
\]
Then
\begin{equation}
3^{k_i}x_{i+1}=4x_i-r_i.
\label{eq:cycle-step}
\end{equation}
Let
\begin{equation}
K=\sum_{i=0}^{L-1}k_i.
\end{equation}
Repeated substitution of \eqref{eq:cycle-step} around the entire loop gives
\begin{equation}
\boxed{
 x_0=
 \frac{\displaystyle
 \sum_{j=0}^{L-1} r_j4^{L-1-j}
 3^{\sum_{m=0}^{j-1}k_m}}
 {4^L-3^K}
 }
\label{eq:closedcycle}
\end{equation}
with the empty exponent sum interpreted as $0$.

Equation~\eqref{eq:closedcycle} gives a necessary exact divisibility condition: for a prescribed residue/valuation word, $4^L-3^K$ must divide the numerator, and the resulting integer states must reproduce the assumed residues and valuations.

\subsection{Sign condition and the critical ratio}
The numerator of \eqref{eq:closedcycle} is positive.  Hence a positive cycle requires
\begin{equation}
4^L-3^K>0
\quad\Longleftrightarrow\quad
\frac KL<\frac{\log4}{\log3}
\approx1.261860.
\label{eq:positivecondition}
\end{equation}
A negative cycle instead requires
\begin{equation}
4^L-3^K<0
\quad\Longleftrightarrow\quad
\frac KL>\frac{\log4}{\log3}.
\label{eq:negativecondition}
\end{equation}
Thus the rational approximation of $\log4/\log3$ by $K/L$ is built directly into the existence problem for cycles.

\section{Algebraic reconstruction of the known cycles}
\subsection{The fixed points}
For a one-cycle, $L=1$ and $K=1$.  With $r_0=1$,
\[
x_0=\frac{1}{4-3}=1.
\]
With $r_0=2$,
\[
x_0=\frac{2}{4-3}=2.
\]
Thus the two positive fixed points arise immediately from the cycle equation.

\subsection{The positive four-cycle}
For \eqref{eq:fourcycle},
\[
L=4,\qquad
(r_0,r_1,r_2,r_3)=(1,2,2,2),
\qquad
(k_0,k_1,k_2,k_3)=(1,1,1,2),
\]
so $K=5$.  The denominator is
\begin{equation}
4^4-3^5=256-243=13.
\end{equation}
The numerator is
\[
1\cdot4^3
+2\cdot4^2\cdot3
+2\cdot4\cdot3^2
+2\cdot3^3
=64+96+72+54=286.
\]
Hence
\begin{equation}
x_0=\frac{286}{13}=22,
\end{equation}
and direct iteration reconstructs
$22\to29\to38\to50\to22$.

The structural point is the unusually small denominator produced by
\begin{equation}
4^4=256\approx243=3^5.
\end{equation}
Equivalently, $K/L=5/4=1.25$ is close to the critical ratio $\log4/\log3\approx1.26186$.  This Diophantine near-balance makes the four-cycle arithmetically possible.

\subsection{Negative cycles as an algebraic extension}
Although the exhaustive $10^{11}$ audit in this paper concerns positive starting values, the same cycle equation also reconstructs the negative cycles noted in the earlier analysis.

For the two-cycle $-1\to-2\to-1$, starting at $-1$ gives
\[
(r_0,r_1)=(2,1),\qquad(k_0,k_1)=(1,2),
\qquad L=2,\ K=3.
\]
Then
\[
4^2-3^3=16-27=-11,
\qquad
2\cdot4+1\cdot3=11,
\]
and therefore $x_0=11/(-11)=-1$.

For the seven-cycle
\begin{equation}
-7\to-10\to-14\to-19\to-26\to-35\to-47\to-7,
\end{equation}
we have
\[
L=7,\quad K=9,
\]
with residue vector $(2,2,1,2,1,1,1)$ and valuation vector $(1,1,1,1,1,1,3)$.  The denominator is
\[
4^7-3^9=16384-19683=-3299,
\]
while the numerator in \eqref{eq:closedcycle} is $23093$.  Thus
\[
x_0=\frac{23093}{-3299}=-7,
\]
which reconstructs the full seven-cycle exactly.  These negative examples emphasize that the sign of $4^L-3^K$ cleanly separates the arithmetic regimes of positive and negative cycles.

\section{Computational methodology}
The finite verification through $10^{11}$ is the union of two disjoint exhaustive computations.  The first component is the previously archived scan from the base article~\cite{banazadeh2026base}, covering
\[
1\le x\le10^9,\qquad 3\nmid x,
\]
with exactly $666,666,667$ admissible starts.  The new V2 extension covers
\[
10^9<x\le10^{11},\qquad 3\nmid x,
\]
with exactly $66,000,000,000$ admissible starts.  Their union is therefore the complete finite domain
\[
1\le x\le10^{11},\qquad3\nmid x,
\]
whose exact cardinality is
\begin{equation}
10^{11}-\left\lfloor\frac{10^{11}}{3}\right\rfloor
=66,666,666,667.
\end{equation}

The V2 extension scanner is implemented in C and iterates the exact accelerated map~\eqref{eq:T}.  Orbit arithmetic uses \texttt{unsigned \_\_int128}; a 64-bit fast path is used only when the affine transition is provably safe in \texttt{uint64\_t}, with automatic promotion to 128-bit arithmetic otherwise.  The scanner checks the two fixed points and the known four-cycle before each transition, and a repeated state outside those known attractors is detected by an exact constant-memory cycle detector and recorded as an unknown cycle.  There is no iteration cap in the V2 production loop: a trajectory terminates only at a known attractor, a newly detected cycle, or an arithmetic boundary event.

The extension interval was divided into $990$ raw-integer chunks of size $100,000,000$ and scanned with four POSIX worker threads.  Each completed chunk is written to a \texttt{.done} checkpoint by atomic rename.  Re-running the scanner skips completed chunks and reconstructs the session totals from the saved records.  Every checkpoint stores the attractor counts, unknown-cycle and overflow counts, total reduced steps, total removed $3$-adic valuation, and the chunk maxima for transient length, accumulated valuation, and fully reduced peak.

\section{V2 validation and reproducibility checks}
Before the $10^{11}$ extension, the V2 implementation was run on
\[
1\le x\le10^6,
\]
using ten raw-integer chunks of size $100,000$.  It exactly reproduces the validation statistics in the base article: $666,667$ admissible starts, with the basin counts shown below.
\begin{center}
\begin{tabular}{lrr}
\toprule
Attractor & Count & Percentage\\
\midrule
$\{1\}$ & 66,382 & 9.957295\%\\
$\{2\}$ & 431,888 & 64.783168\%\\
$22\to29\to38\to50\to22$ & 168,397 & 25.259537\%\\
\bottomrule
\end{tabular}
\end{center}
The same validation recovers the $249$-iteration transient record from $969,812$, the accumulated valuation record $326$ from the same start, and the fully reduced peak $68,682,968,834$ from $491,732$.

A separate V2 benchmark run on the first production chunk
\[
1,000,000,001\le x\le1,100,000,000
\]
reproduced the saved production checkpoint for that interval in all integer-valued statistics.  This duplicate calculation provides an additional check on the production executable and checkpoint format.

\section{New extension scan: $10^9<x\le10^{11}$}
The $990$ production checkpoints cover the extension interval contiguously from $1,000,000,001$ through $100,000,000,000$.  Summing their saved statistics gives
\begin{equation}
\boxed{\text{tested}=66,000,000,000,\quad
\text{unknown}=0,\quad
\text{overflow}=0.}
\end{equation}
No terminal positive cycle other than the two fixed points and the four-cycle was found.

\begin{table}[ht]
\centering
\caption{Attractor classification in the new interval $10^9<x\le10^{11}$, $3\nmid x$.}
\label{tab:extension}
\begin{tabular}{lrr}
\toprule
Attractor & Number of starts & Percentage\\
\midrule
$\{1\}$ & 6,536,905,587 & 9.904402\%\\
$\{2\}$ & 42,779,024,148 & 64.816703\%\\
$22\to29\to38\to50\to22$ & 16,684,070,265 & 25.278894\%\\
\midrule
Total & 66,000,000,000 & 100.000000\%\\
\bottomrule
\end{tabular}
\end{table}

The exact count identity
\[
6,536,905,587+42,779,024,148+16,684,070,265=66,000,000,000
\]
accounts for every admissible start in the extension interval.  The same checkpoints record a total of $5,533,148,453,392$ reduced-map iterations and a total accumulated $3$-adic valuation of $8,362,466,366,261$ over the extension scan.

The extension produces new finite-domain records.  The longest reduced-map transient is
\begin{equation}
\boxed{574\text{ iterations from }x_0=97,312,923,676,}
\end{equation}
and the accumulated number of removed factors of $3$ from the same start is
\begin{equation}
\boxed{744.}
\end{equation}
The largest fully reduced state recorded in the extension is
\begin{equation}
\boxed{580,315,803,007,928,858,285,}
\end{equation}
reached on the trajectory beginning at
\begin{equation}
\boxed{x_0=52,991,345,843.}
\end{equation}
All three records exceed the corresponding records in the base $10^9$ computation.

\section{Combined exhaustive verification through $10^{11}$}
The base interval and the new extension are disjoint and together exhaust every positive admissible start through $10^{11}$.  Adding the base counts~\cite{banazadeh2026base} to Table~\ref{tab:extension} gives the complete classification.

\begin{table}[ht]
\centering
\caption{Complete attractor classification for all $1\le x\le10^{11}$ with $3\nmid x$.}
\label{tab:1e11}
\begin{tabular}{lrr}
\toprule
Attractor & Number of starts & Percentage\\
\midrule
$\{1\}$ & 6,602,925,203 & 9.904388\%\\
$\{2\}$ & 43,211,124,603 & 64.816687\%\\
$22\to29\to38\to50\to22$ & 16,852,616,861 & 25.278925\%\\
\midrule
Total & 66,666,666,667 & 100.000000\%\\
\bottomrule
\end{tabular}
\end{table}

The combined count identity
\[
6,602,925,203+43,211,124,603+16,852,616,861
=66,666,666,667
\]
is exact.  Since both component scans report zero unknown cycles and zero arithmetic overflows, the same is true of the combined finite verification.  Every tested start through $10^{11}$ enters one of the two fixed points or the positive four-cycle.


The V2 scanner was also applied to the complete base interval $1\le x\le10^9$ as a supplementary aggregate verification.  It reproduces the archived base basin counts and all three published base extrema exactly, while additionally recording
\[
\text{base total reduced steps}=44,033,690,920,
\qquad
\text{base total }\vthree=66,685,065,543.
\]
The base total-step value independently agrees with the exact integer recovered from the archived chunk means in the original $10^9$ computational package.  Adding these base aggregates to the disjoint V2 extension totals gives the following full-domain statistics.

\begin{table}[ht]
\centering
\small
\caption{Aggregate trajectory statistics for the two disjoint computational intervals and their union.}
\label{tab:aggregate}
\begin{tabular}{lrrr}
\toprule
Quantity & Base $x\le10^9$ & Extension & Combined $x\le10^{11}$\\
\midrule
Admissible starts & 666,666,667 & 66,000,000,000 & 66,666,666,667\\
Total reduced iterations & 44,033,690,920 & 5,533,148,453,392 & 5,577,182,144,312\\
Mean reduced iterations/start & 66.050536347 & 83.835582627 & 83.657732164\\
Total removed $3$-adic factors & 66,685,065,543 & 8,362,466,366,261 & 8,429,151,431,804\\
Mean removed factors/start & 100.027598264 & 126.704035852 & 126.437271476\\
Unknown cycles & 0 & 0 & 0\\
Arithmetic overflows & 0 & 0 & 0\\
\bottomrule
\end{tabular}
\end{table}

For transparency, the arithmetic decomposition of the three basin counts is
\begin{center}
\small
\begin{tabular}{lrrr}
\toprule
Attractor & Base $x\le10^9$ & Extension & Combined\\
\midrule
$\{1\}$ & 66,019,616 & 6,536,905,587 & 6,602,925,203\\
$\{2\}$ & 432,100,455 & 42,779,024,148 & 43,211,124,603\\
Four-cycle & 168,546,596 & 16,684,070,265 & 16,852,616,861\\
\midrule
Total & 666,666,667 & 66,000,000,000 & 66,666,666,667\\
\bottomrule
\end{tabular}
\end{center}

\section{Comparison across verified scales}
The basin proportions remain numerically stable as the verified domain expands from $10^6$ to $10^9$ and then to $10^{11}$.
\begin{table}[ht]
\centering
\caption{Observed positive basin proportions across three exhaustive cutoffs.}
\label{tab:compare}
\begin{tabular}{lrrr}
\toprule
Attractor & $x\le10^6$ & $x\le10^9$ & $x\le10^{11}$\\
\midrule
$\{1\}$ & 9.957295\% & 9.902942\% & 9.904388\%\\
$\{2\}$ & 64.783168\% & 64.815068\% & 64.816687\%\\
Four-cycle & 25.259537\% & 25.281989\% & 25.278925\%\\
\bottomrule
\end{tabular}
\end{table}
The proportions move only slightly across these ranges.  This is compatible with the existence of limiting basin densities, but the finite experiments neither prove that such limits exist nor determine exact limiting values.

The extremal orbit statistics, by contrast, continue to grow sharply.
\begin{table}[ht]
\centering
\caption{Growth of selected finite-domain records.}
\label{tab:records}
\small
\begin{tabular}{lrrr}
\toprule
Cutoff & Max. reduced steps & Max. accumulated $\vthree$ & Max. reduced peak\\
\midrule
$10^6$ & 249 & 326 & 68,682,968,834\\
$10^9$ & 405 & 529 & 164,210,561,845,359,809\\
$10^{11}$ & 574 & 744 & 580,315,803,007,928,858,285\\
\bottomrule
\end{tabular}
\end{table}
Thus stable aggregate basin proportions coexist with increasingly extreme individual excursions.

\section{Interpretation of the cycle structure}
Three mechanisms interact in the dynamics.

First, the valuation distribution favors contraction on the logarithmic average scale.  Second, the exact cycle equation imposes a strong divisibility constraint on every periodic residue--valuation word.  Third, the balance between $4^L$ and $3^K$ controls the size and sign of the cycle denominator.  The positive four-cycle is a particularly transparent example because $4^4-3^5=13$ is unusually small.

This suggests a useful search principle for periodic orbits: rational approximations $K/L$ near $\log4/\log3$ produce small values of $|4^L-3^K|$ and hence arithmetically favorable candidates, but a candidate still has to satisfy the exact numerator divisibility and reproduce its residue and valuation pattern.  The near equality alone is therefore not sufficient for a cycle.

\section{Finite theorem and global conjecture}
The combined computation establishes the following finite statement.

\medskip
\noindent\textbf{Finite computational statement.}
For every $x\in\Sset$ with $1\le x\le10^{11}$, the exact finite audit represented by the base computation~\cite{banazadeh2026base} together with the archived V2 extension reaches one of
\[
\{1\},\qquad\{2\},\qquad
22\to29\to38\to50\to22,
\]
and neither component computation reports an arithmetic overflow or an additional terminal cycle.
\medskip

The numerical evidence motivates the same global conjecture as before.

\medskip
\noindent\textbf{Conjecture (positive three-attractor convergence).}
Every positive integer not divisible by $3$, under the map~\eqref{eq:T}, eventually enters one of the two fixed points $1,2$ or the four-cycle $22\to29\to38\to50\to22$.
\medskip

The exhaustive finite verification through $10^{11}$ is substantially stronger evidence than the earlier $10^9$ audit, but it does not prove the conjecture.  A proof must exclude every additional positive cycle at all larger scales and rule out unbounded trajectories.

\section{What remains for a proof?}
A complete proof would still need to overcome at least three distinct obstacles.
\begin{enumerate}[leftmargin=2em]
\item The geometric distribution~\eqref{eq:geom} is a probabilistic model.  A deterministic proof needs sufficient control of the actual valuation frequencies along every orbit, or an argument that does not depend on statistical independence.
\item The cycle equation~\eqref{eq:closedcycle} must be controlled for arbitrarily large $L$ and $K$.  Finite searches over lengths cannot exclude a cycle with much larger parameters.
\item Large temporary excursions must be handled uniformly.  The observed reduced peak above $5.8\times10^{20}$ from a start below $10^{11}$ demonstrates that early descent or small local growth cannot serve as a global argument.
\end{enumerate}
These limitations separate the rigorous finite computation from the still-open all-integers statement.

\section{Reproducibility}
The computational archive accompanying this manuscript contains only V2 materials: the V2 C source and corresponding executable, ten validation checkpoints, one benchmark checkpoint, all $990$ production checkpoints for the $10^9<x\le10^{11}$ extension, ten supplementary V2 checkpoints covering the base interval $1\le x\le10^9$, a combined aggregate summary, and the final LaTeX source of this manuscript.  The PDF is distributed separately.  No V1, V3, V4, or V5 scanner files are included.

A typical V2 compilation command is
\begin{verbatim}
clang -O3 -march=native -pthread scan_ternary_minus_1e11_v2.c \
  -o scan_ternary_minus_1e11_v2
\end{verbatim}
The production extension scan is represented by
\begin{verbatim}
./scan_ternary_minus_1e11_v2 1000000001 100000000000 \
  100000000 4 results_1e11_chunks
\end{verbatim}
and the validation run through $10^6$ by
\begin{verbatim}
./scan_ternary_minus_1e11_v2 1 1000000 100000 4 \
  test_v2_chunks
\end{verbatim}
The supplementary V2 base verification is reproduced by
\begin{verbatim}
./scan_ternary_minus_1e11_v2 1 1000000000 100000000 4 \
  base_1e9_v2_verification
\end{verbatim}

The archived chunk files themselves are the machine-readable audit trail.  Independent aggregation of the $990$ production records should recover the extension counts in Table~\ref{tab:extension}, zero unknown cycles, zero overflow events, the $574$-iteration transient record, the accumulated valuation record $744$, and the maximum fully reduced state reported above.  Aggregating the ten supplementary base checkpoints reproduces the published $10^9$ basin counts and extrema and yields the base totals in Table~\ref{tab:aggregate}; summing the two disjoint intervals yields the combined totals $5,577,182,144,312$ reduced iterations and $8,429,151,431,804$ removed factors of $3$.

\section{Conclusion}
The ternary $(4k-1(2))/3$ Collatz-type system studied here is most precisely represented by the accelerated map
\[
\Tmap(x)=\frac{4x-(x\bmod3)}{3^{\vthree(4x-(x\bmod3))}},
\qquad3\nmid x.
\]
Its local rule is elementary, but the resulting arithmetic dynamics combine residue branching, variable $3$-adic compression, large transient excursions, and nontrivial periodic structure.

The theoretical analysis is unchanged in character from the base study: the geometric valuation model gives a negative average logarithmic drift and the exact cycle equation converts periodicity into a Diophantine divisibility problem controlled by the balance of $4^L$ and $3^K$.  The new contribution is computational scale.  Combining the archived $10^9$ audit with the disjoint V2 extension classifies all $66,666,666,667$ admissible positive starts through $10^{11}$.  Across this complete domain the trajectories contain $5,577,182,144,312$ reduced iterations and $8,429,151,431,804$ removed factors of $3$.  No additional positive attractor, unresolved cycle, or overflow is reported, while the record transient, accumulated valuation, and reduced peak all increase substantially.  The resulting three-attractor convergence conjecture nevertheless remains open beyond the finite verified range.

\section*{Data and code availability}
The Zenodo deposit for this $10^{11}$ study is intended to contain three deliverables: the PDF manuscript, the standalone LaTeX source, and a V2-only computational archive containing the scanner, validation and benchmark checkpoints, the complete $990$-chunk extension audit trail, the ten-checkpoint supplementary V2 base verification, the combined aggregate summary, and the same LaTeX source.  The earlier $10^9$ paper and its finite results remain independently citable through its Zenodo DOI below.

\begin{thebibliography}{9}\sloppy
\bibitem{lagarias1985}
J. C. Lagarias,
``The $3x+1$ problem and its generalizations,''
\emph{The American Mathematical Monthly}, 92(1), 3--23, 1985.

\bibitem{carnielli2015}
W. Carnielli,
``Some Natural Generalizations of the Collatz Problem,''
\emph{Applied Mathematics E-Notes}, 15, 207--212, 2015.

\bibitem{banazadeh2026base}
F. Banazadeh,
\emph{A Ternary $(4k-1(2))/3$ Collatz-Type Map: Complete 3-Adic Reduction, Algebraic Cycle Analysis, and Exhaustive Computational Verification up to $10^9$},
Zenodo, 2026.\\
\url{https://doi.org/10.5281/zenodo.22662980}

\end{thebibliography}

\end{document}
